\chapter*{Appendix - Glossary}
\addcontentsline{toc}{chapter}{Appendix - Glossary}

\paragraph{Absolute Address}
An address in the physical memory range from zero to the maximum physical memory size. This range is directly mapped to an equivalent virtual address range.

\paragraph{Address Space}
A logical domain of addresses. The CPU may define multiple address spaces, such as \textit{main memory space} and \textit{I/O space}. Each space defines an independent address range and access semantics.

\paragraph{Callee Save Registers}
Registers whose values must be preserved across a function call by the callee (the function being called).  
If the callee uses any of these registers, it must save their original values on entry (typically on the stack) and restore them before returning.  
This convention allows callers to assume the registers retain their values after the call.

\paragraph{Caller Save Registers}
Registers that may be overwritten by a called function.  
The caller must save these registers (if it needs their values later) before invoking another function.  
This convention allows the callee to use these registers freely without preserving their previous contents.

\paragraph{ELF Section}
A named subdivision in an ELF file used by the linker and compiler (e.g., \texttt{.text}, \texttt{.data}, \texttt{.bss}). Sections are grouped by the linker into ELF segments.

\paragraph{ELF Segment}
A loadable portion of an ELF file (defined by a program header). Each ELF segment is mapped into one or more regions of an address space during program loading.

\paragraph{Function Pointer}
A reference to a function descriptor.

\paragraph{Global Data Area}
All modules compiled for a load module allocate their global data in this area. The DP register points to the beginning of the area.

\paragraph{Global Data Pointer (DP)}
The address of the global data area. The address is normally kept in the DP register. Each module has a DP area in the global data area. Load modules also use this register to access the linkage table.

\paragraph{IA-relative Addressing}
Code that uses the current instruction address (IA) as the base for branches and for accessing constant data.

\paragraph{Linkage Table}
A table of address descriptors for functions external to the load module. One type of entry stores the address of an external function. Another type stores a reference to the external load module's global data area.

\paragraph{Load Module}
An executable unit of code produced by the assembler or compiler from one or several modules.

\paragraph{Mapped Region}
A runtime mapping that associates an ELF segment or memory source with a region in an address space. The emulator or loader tracks mapped regions to manage memory access and protection.

\paragraph{Object}
A unit of storage within a region. Objects correspond to program-visible entities or runtime allocations, such as global variables, heap blocks, or stack frames.  
Examples: \textit{global object}, \textit{heap object}, \textit{stack object}.

\paragraph{Offset}
An unsigned 32-bit value that, together with the region, forms a virtual address.  
Examples: \textit{code offset}, \textit{data offset}.

\paragraph{Region}
A contiguous range of addresses within the virtual address space, reserved for a particular purpose. A region represents the upper 20 bits of the 52-bit virtual address range. 
Examples: \textit{code region}, \textit{data region}, \textit{stack region}.

\paragraph{Virtual Address}
A virtual address consists of the region ID and the region offset. Together they form the architected virtual address space of 52 bits.
